Quantum Boltzmann Machines: Theory and ImplementationOA







Preface/Introduction





Quantum Machine Learning (QML) combines the power of quantum computing
with machine learning models. A particularly intriguing model in this
landscape is the Quantum Boltzmann Machine (QBM), which extends the
classical Boltzmann machine (a probabilistic neural network) into the
quantum domain. QBMs promise richer representations by leveraging
quantum superposition and entanglement, potentially capturing
correlations that classical models cannot . In these notes, we review
classical Boltzmann machines and restricted Boltzmann machines (RBMs),
introduce QBMs and their restricted variant (RQBM), discuss training
methods, and illustrate practical implementation using
PennyLane. Mathematical derivations are provided, along with clear
diagrams and quantum code examples. Each chapter ends with exercises
to deepen understanding.





Table of Contents





Classical Boltzmann Machines and RBMs Review
Introduction to Quantum Boltzmann Machines (QBMs)
Restricted Quantum Boltzmann Machines (RQBM)
Training Quantum Boltzmann Machines
Practical Implementation with PennyLane
Applications and Examples
References
Index






Classical Boltzmann Machines and RBMs Review





A Boltzmann machine (BM) is an undirected probabilistic graphical model of binary units (neurons) with pairwise interactions . Each unit $i$ has a binary state $x_i\in{0,1}$ (or $\pm1$) and an associated bias $b_i$, and each pair of units $(i,j)$ may have a symmetric interaction weight $w_{ij}$. The energy of a state $x=(x_1,\dots,x_N)$ is defined as

$$

E(x; b,w) = -\sum_{i<j} w_{ij},x_i x_j - \sum_{i} b_i,x_i,,

$$

where we adopt the physics convention (hence the minus signs) so that lower energy states are more probable.  The BM defines a Gibbs (Boltzmann) distribution over states:

$$

p(x; b,w) ;=; \frac{1}{Z(b,w)} \exp(-E(x; b,w)),

\qquad Z(b,w)=\sum_{x\in{0,1}^N} \exp(-E(x; b,w)),,

$$

where $Z$ is the partition function (normalization) .  In other words, $p(x)\propto e^{-E(x)}$ and $Z$ sums $e^{-E}$ over all $2^N$ configurations .  A BM is a general energy-based model with hidden (unobserved) and visible (observed) units. In learning, the goal is to adjust $(b,w)$ so that the model distribution matches the empirical data distribution.



A Restricted Boltzmann Machine (RBM) is a special case of BM with a bipartite structure .  Units are divided into a visible layer $v\in{0,1}^n$ (observed data) and a hidden layer $h\in{0,1}^m$, and no connections exist among units within the same layer. This restriction simplifies many calculations. The energy of an RBM state is

$$

E(v,h; W,b,c) = -\sum_{i,j} W_{ij},h_j v_i - \sum_i b_i v_i - \sum_j c_j h_j,

$$

where $W_{ij}$ couples visible unit $i$ to hidden $j$, $b_i$ are visible biases and $c_j$ hidden biases. The joint probability is

$$

p(v,h) = \frac{1}{Z} e^{-E(v,h)},

$$

and the marginal probability over visible states is

$$

p(v) = \sum_{h} p(v,h) = \frac{1}{Z} \sum_{h\in{0,1}^m} \exp\bigl(\sum_{i,j} W_{ij},h_j v_i + \sum_i b_i v_i + \sum_j c_j h_j\bigr)!.

$$

Equivalently, one can show :

$$

p(v) = \frac{1}{Z} \exp\Bigl(\sum_i b_i v_i\Bigr)\prod_{j=1}^m \bigl(1 + e^{c_j + \sum_i W_{ij}v_i}\bigr).

$$

Figure below illustrates an RBM network architecture.



Figure: Restricted Boltzmann Machine (RBM). Visible units $v_i$ (green) are connected to hidden units $h_j$ (blue), with weights $W_{ij}$. No visible-visible or hidden-hidden edges exist, enabling conditional independence and efficient sampling .



Learning in a BM/RBM: Training adjusts weights so that the model reproduces data statistics. One minimizes the negative log-likelihood $-\sum_{data}\ln p(v)$ of the data. The gradient of this loss w.r.t.\ a parameter (e.g.\ a weight or bias) involves the difference of correlations between the data and the model. For example, for a general BM the update rule (stochastic gradient descent) is :

$$

\delta W_{ij} ;\propto; \langle x_i x_j\rangle_{\text{data}} ;-; \langle x_i x_j\rangle_{\text{model}}!, \qquad

\delta b_{i} ;\propto; \langle x_i \rangle_{\text{data}} ;-; \langle x_i \rangle_{\text{model}}!,

$$

where $\langle\cdot\rangle_{\text{data}}$ are averages over data samples and $\langle\cdot\rangle_{\text{model}}$ are expectations under the model distribution. Computing $\langle\cdot\rangle_{\text{model}}$ requires (in general) summing over all states or performing Gibbs sampling, which can be expensive. In practice for RBMs one often uses Contrastive Divergence (CD) or similar approximations to efficiently approximate these expectations.



RBM Sampling and CD: In an RBM, given visible $v$, the hidden units are conditionally independent:

$$

p(h_j=1\mid v) = \sigma!\bigl(c_j + \sum_i W_{ij} v_i\bigr)!,

$$

where $\sigma(x) = 1/(1+e^{-x})$ is the logistic function. Similarly, given hidden $h$, the visibles are independent. This bipartite structure allows one Gibbs step (sampling $h$ from $p(h|v)$, then $v$ from $p(v|h)$) efficiently. Contrastive Divergence training uses just a few Gibbs steps from the data to approximate the model averages .



Key properties: Classical Boltzmann machines form an exponential family of probability distributions. They can approximate any distribution arbitrarily well if sufficiently large (universal approximator), but training is generally hard (NP-hard) due to the partition function $Z$. RBMs, due to their structure, are easier to train and have been widely used in unsupervised learning (e.g.\ as building blocks of deep belief networks ).



Exercises (Chapter 2): (1) Compute partition function: For a 2-bit fully connected BM with weights $w_{12}=0.5$, biases $b_1=b_2=0$, list the energy of each state and compute $Z$ and the probability of state $(1,1)$. (2) RBM conditional probabilities: Write the expressions for $p(h_j=1\mid v)$ and $p(v_i=1\mid h)$ in terms of the RBM parameters. (3) Contrastive Divergence: Outline the CD-1 weight update rule in words and explain why it is an approximation to maximum likelihood learning.





Introduction to Quantum Boltzmann Machines (QBMs)





A Quantum Boltzmann Machine (QBM) extends a classical BM by replacing each binary unit with a qubit and generalizing the energy to a quantum Hamiltonian .  Concretely, consider a system of $N$ qubits. A convenient choice is the transverse-field Ising model (TFIM) Hamiltonian:

$$

H = -\sum_{a=1}^N \Gamma_a,\sigma^x_a ;-;\sum_{a=1}^N b_a,\sigma^z_a ;-;\sum_{a<b} u_{ab},\sigma^z_a \sigma^z_b,

$$

where $\sigma_a^x,\sigma_a^z$ are the Pauli matrices on qubit $a$, $\Gamma_a$ is a “transverse” field, $b_a$ are local fields (biases), and $u_{ab}$ are interaction strengths .  Here each qubit can be interpreted analogously to a classical unit, but the $\sigma^x$ term induces quantum superposition.  One can also include $\sigma^y$ or more complex terms, but we focus on this common ansatz.



The quantum model assigns a density matrix (mixed state) to the system via the Gibbs (thermal) state at some inverse temperature $\beta=1/T$:

$$

\rho ;=; \frac{e^{-\beta H}}{\Tr(e^{-\beta H})},.

$$

Measuring $\rho$ in the computational ($\sigma^z$) basis yields a probability distribution over bit-strings $z\in{0,1}^N$:

$$

p(z) = \bra{z} \rho \ket{z} ;=; \frac{\bra{z} e^{-\beta H} \ket{z}}{\Tr(e^{-\beta H})}.

$$

This model is a quantum generalization of the Boltzmann distribution: if $H$ were diagonal in the $z$-basis (i.e.\ $\Gamma_a=0$), then $\bra{z} e^{-\beta H} \ket{z} = e^{-\beta E(z)}$ and one recovers a classical BM. But with nonzero $\Gamma$, $H$ has off-diagonal terms, and $\rho$ generally has entanglement and non-commuting contributions.



One advantage of QBMs is expressivity: quantum Hamiltonians can capture correlations that classical ones cannot. For instance, the non-commuting $\sigma^x$ terms allow exploring multiple states in superposition. Studies suggest that even when modeling classical data, QBMs can outperform classical BMs in capturing complex patterns .  Askarzadeh et al. note that a QBM can include “more general non-commuting terms,” making it strictly more expressive .  In practice, one often treats the visible data as classical: one may fix (clamp) the “visible” qubits to data values and only allow quantum sampling over hidden qubits, similar to RBMs. But fundamentally, a QBM defines a quantum Boltzmann distribution as above .



In summary, a QBM is a physics-inspired generative model: its parameters $(\Gamma, b, u)$ define a Hamiltonian, and sampling is done by preparing the Gibbs state of that Hamiltonian. By training these parameters, the QBM is intended to learn to generate data similar to the training set.  As with classical BMs, one can consider visible and hidden qubits: visible qubits may be initialized to classical data (e.g.\ via computational-basis initialization), while hidden qubits are free to evolve quantumly.  The model is trained so that measurements on the visible qubits (after tracing out or measuring hidden qubits) reproduce the data distribution .



Why quantum?  Intuitively, using a quantum system can allow more compact representations.  For example, Gurvits and others have shown that quantum statistics can represent certain probability distributions more efficiently than classical spins.  Moreover, upcoming quantum hardware (e.g.\ quantum annealers) may naturally implement sampling from such quantum Boltzmann distributions.  While the ultimate advantage is still under research, the idea is that QBMs could be components of future quantum machine learning systems, possibly providing speedups or capturing quantum correlations in quantum data .



Representative fact:  Tan et al. (2023) remark that a QBM is “a physically motivated quantum neural network” for generative/discriminative tasks . Similarly, Amin et al. (2018) introduce a QBM learning algorithm and demonstrate via simulations that “the machine exploits its quantum nature to mimic data sets” .



Exercises (Chapter 3): (1) QBM Hamiltonian: Write the Hamiltonian $H$ for a QBM with two visible qubits and one hidden qubit, using bias parameters $(b_1,b_2)$ for visibles, $(c_1)$ for hidden, weights $W_{1},W_{2}$ between visible 1 and hidden, etc., plus transverse fields $\Gamma_i$. (2) Quantum vs Classical: In the limit $\Gamma_a\to 0$, show that the QBM reduces to a classical BM. What role does $\Gamma$ play? (3) Density matrix: Explain why the thermal state $\rho=e^{-\beta H}/Z$ is a density matrix and how measuring in the computational basis yields a probability distribution over bit-strings.





Restricted Quantum Boltzmann Machines (RQBM)





As in the classical case, one can restrict a QBM’s connectivity to simplify training. A Restricted QBM (RQBM) has a bipartite graph structure: visible qubits connect only to hidden qubits, with no interactions among visibles or among hiddens.  The Hamiltonian then has couplings $u_{v,h}$ only between a visible qubit $v$ and a hidden qubit $h$.  Formally, for $n$ visible and $m$ hidden qubits, one can write:

$$

H = -\sum_{i=1}^n b_i,\sigma^z_{v_i} - \sum_{j=1}^m c_j,\sigma^z_{h_j} - \sum_{i,j} W_{ij},\sigma^z_{v_i}\sigma^z_{h_j} ;-; \sum_{j=1}^m \Gamma_j,\sigma^x_{h_j},

$$

where we may allow transverse fields only on the hidden qubits (visible qubits are clamped to data and treated classically). This is a quantum analog of the classical RBM: visible units (qubits) have no $\sigma^z$–$\sigma^z$ interactions among themselves, only with hidden qubits.



Such an architecture can make training easier.  Wiebe and Wossnig (2019) specifically consider a class of RQBMs where the hidden-qubit Hamiltonians commute, allowing certain variational training techniques . They point out that classical pre-training or other bounds can be used in the restricted case.  In reinforcement learning applications, Amin et al. have applied an RQBM to learn policies by clamping visible inputs (state-action pairs) and letting hidden qubits represent the “value” function .



In general, the RQBM still defines a quantum Boltzmann distribution: the Gibbs state of the above Hamiltonian yields

$$

\rho = \frac{e^{-\beta H}}{\Tr e^{-\beta H}}, \qquad

P(v) = \Tr_{h}\bigl[\bra{v} \rho \ket{v}\bigr],

$$

where $\Tr_h$ denotes tracing over hidden qubits. Training then adjusts ${W,b,c,\Gamma}$ so that $P(v)$ matches the data (where $v$ runs over visible bit-strings). By enforcing the bipartite connectivity, one hopes that sampling and gradient estimates can be done more efficiently, akin to classical RBM training.



Analogy: The RQBM retains the spirit of a classical RBM but with quantum dynamics on the hidden layer. One can view it as a hybrid: the visible qubits are often “classical” (fixed in the computational basis for each data point), while hidden qubits evolve under a transverse-field Ising Hamiltonian. This structure has been used both theoretically and experimentally to realize QBMs on hardware like quantum annealers .



Exercises (Chapter 4): (1) Graph structure: Draw an RQBM with 2 visible and 2 hidden qubits, labeling the Hamiltonian terms (biases on all qubits and weights only between visible–hidden pairs). (2) Commuting hidden Hamiltonians: Why might one require hidden-qubit Hamiltonians to commute for certain training algorithms? Discuss the case where $\Gamma$ terms are equal on all hidden qubits. (3) Connection to quantum annealing: Explain how an RQBM Hamiltonian is similar to the Hamiltonians used in quantum annealing devices (e.g.\ D-Wave), and why annealers naturally sample from (or approximate) such quantum Boltzmann distributions.





Training Quantum Boltzmann Machines





Training a QBM means optimizing its parameters so that the model distribution $\rho$ fits the data distribution. A natural objective is the quantum relative entropy (quantum Kullback–Leibler divergence) between the data ensemble and the model’s Gibbs state. Equivalently, one minimizes the negative log-likelihood of the data under the QBM.  The gradient of this loss with respect to a parameter $\theta$ typically involves terms like $\partial_\theta \Tr(\rho \ln \rho_D)$, which translate to differences of expectation values. For a QBM with Hamiltonian $H(\theta)$, it can be shown that the gradient of the log-likelihood leads to updates of the form :

$$

\delta \theta ;\propto; \langle O_\theta \rangle_{\text{data}} ;-; \langle O_\theta \rangle_{\text{model}},

$$

where $O_\theta = \partial_\theta H$ is the operator conjugate to $\theta$, and the expectations are taken over the data-clamped and model Gibbs states respectively. For example, for a weight $W_{ij}$ the gradient involves $\langle \sigma^z_{v_i}\sigma^z_{h_j}\rangle$, and for a bias $b_i$ it involves $\langle\sigma^z_{v_i}\rangle$.



However, unlike the classical case, computing the model expectations $\langle O_\theta\rangle_{\text{model}}=\Tr(\rho,O_\theta)$ requires sampling from a quantum thermal state $\rho=e^{-\beta H}/Z$.  Preparing such a Gibbs state on a quantum computer is generally hard (in fact NP-hard in worst case) .  As Tan et al. note, training requires many such samples since each gradient estimate demands potentially exponential resources .



Thus, training QBM is challenging. Several approaches exist:



Quantum variational training: One can use a variational quantum circuit to approximate the Gibbs state. For example, Quantum Imaginary Time Evolution (QITE) or a layered ansatz can approximate $e^{-\beta H}\ket{+}^{\otimes n}$. The parameters of this circuit are then optimized so that its output matches data statistics.  PennyLane and other frameworks can compute gradients via parameter-shift rules for such circuits. (This approach requires a fault-tolerant device or clever heuristics.)
Bounds and approximations: Amin et al. (2018) circumvented the non-commutativity issue by introducing a bound on quantum probabilities, allowing approximate sampling . In their method, they derive a lower bound to the log-likelihood that can be sampled efficiently. They demonstrate training with and without this bound, using exact diagonalization on small systems.  The essence is to replace $e^{-\beta H}$ with an operator that is easier to sample from, controlling the error.
Restricted/commuting methods: Wiebe & Wossnig (2019) developed methods for Restricted QBMs with commuting hidden Hamiltonians. By ensuring that the hidden-qubit Hamiltonian terms commute, one can derive a variational upper bound on the quantum relative entropy and optimize it classically .  In this restricted case, the problem simplifies, and gradient evaluation can be done using linear combinations of unitaries or by sampling from simpler Hamiltonians.
Hybrid annealing: One proposal is to use quantum annealers themselves as Boltzmann samplers. If one encodes the QBM Hamiltonian (with quantum fluctuations) onto a quantum annealer, then running it at finite temperature might (approximately) sample from the needed quantum Gibbs distribution. Some works explore training annealer qubits to mimic a QBM through reverse annealing or similar techniques .
Coreset and data reduction: Recently, Tan et al. (2023) suggested using coresets (small representative data subsets) to reduce training cost . The idea is to compress the training set into a weighted coreset so that the QBM only needs to match a small number of effective data points, reducing the number of Gibbs-state preparations needed.




Overall, training a QBM often involves computing gradients of the form

$$

\partial_\theta \bigl[-\log \Tr(e^{-\beta H})\bigr] ;=; \langle \partial_\theta H \rangle_{\text{model}}!,

$$

which is then set against empirical averages.  In practice, one may replace exact gradients with approximate or stochastic estimates.  Note that when $\Gamma=0$, the model is classical and one recovers standard BM gradients .



Contrast with classical training:  For classical RBMs, the Contrastive Divergence (CD) algorithm provides an efficient approximate gradient by a short Gibbs chain . A quantum analog of CD is less straightforward, due to the need to sample a quantum state. Some proposals use one round of (partial) quantum annealing as the “negative phase” to get an estimate of model correlations. These methods are still experimental.



Exercise: Show that minimizing the quantum relative entropy $S(\rho_D|\rho_\theta)$ (where $\rho_D$ is the data density matrix and $\rho_\theta=e^{-\beta H(\theta)}/Z$) is equivalent to maximizing the log-likelihood of the data under the model. Derive the gradient $\partial_\theta S$ to confirm it yields terms of the form $\langle \partial_\theta H \rangle_{\text{data}} - \langle \partial_\theta H \rangle_{\text{model}}$. 





Practical Implementation with PennyLane





To illustrate QBM concepts, we can use PennyLane – a Python library for quantum machine learning – to implement simple generative models. Below we sketch a toy example: a 2-qubit Quantum Circuit Born Machine (QCBM), which generates a probability distribution via a parameterized quantum circuit. While a QCBM is not exactly a QBM (it has no thermal state), it serves to show how to code a quantum generative model.

import pennylane as qml
from pennylane import numpy as np

# Use a 2-qubit simulator
dev = qml.device('default.qubit', wires=2)

# Define a simple 2-qubit QCBM circuit
@qml.qnode(dev)
def circuit(params):
    # params = 4 angles
    qml.RX(params[0], wires=0)
    qml.RY(params[1], wires=1)
    qml.CNOT(wires=[0, 1])
    qml.RX(params[2], wires=0)
    qml.RY(params[3], wires=1)
    return qml.probs(wires=[0,1])  # return probabilities of |00>,|01>,|10>,|11>

# Initialize random parameters
params = np.random.randn(4, requires_grad=True)

# Get the output distribution from the circuit
probs = circuit(params)
print("Probabilities:", probs)
In this code, we create a 2-qubit variational circuit and return the probabilities of each basis state. One could train params to match a target distribution by defining a cost (e.g. KL divergence) between probs and the data distribution, and using PennyLane’s automatic differentiation to update params.  Though not a true QBM (no Hamiltonian or Gibbs state), this demonstrates how one might build and train a small quantum generative model in PennyLane.



To make it more QBM-like, one could encode the classical visible units as fixed inputs. For example, to compute the energy of a QBM Hamiltonian for a given visible bitstring, one might do:

# Example: compute expectation of a simple Hamiltonian for a given state |v>
H = 1.5 * qml.PauliZ(0) + 0.7 * qml.PauliZ(1) + 0.9 * qml.PauliZ(0)@qml.PauliZ(1)

@qml.qnode(dev)
def energy_of_state():
    # Prepare visible qubits in state |v> = |01>, say
    qml.PauliX(wires=1)  # flips qubit 1 to |1>
    # No hidden qubits in this simple example
    # Return expectation of H in this state
    return qml.expval(H)

print("Energy of |01> state:", energy_of_state())
Here, we used a 2-qubit Hamiltonian $H=1.5\sigma_z^0+0.7\sigma_z^1+0.9\sigma_z^0\sigma_z^1$ and prepared the state $|v>=|01>$. The expval(H) call returns $\langle 01|H|01\rangle = -1.5 + 0.7 - 0.9 = -1.7$ (up to sign conventions). This shows how PennyLane can compute energies of basis states, which is a key step in evaluating the QBM energy and probability of data states.



In practice, training a QBM in PennyLane would involve preparing quantum states corresponding to the Gibbs distribution. One could use functions like qml.qaoa.cost.Hamiltonian or custom circuits to approximate $e^{-\beta H}$, and then use PennyLane’s optimizers. Also, PennyLane can interface with PyTorch or TensorFlow, enabling hybrid optimization of quantum circuits with classical parameters (e.g. the Hamiltonian weights).



Summary of steps in PennyLane implementation:



Define a quantum device (qml.device) with enough wires for visibles+hidden.
Construct a parametric quantum circuit (ansatz) that depends on trainable parameters (e.g. angles in rotations and entangling gates).
If modeling an RQBM, one might clamp visible qubits (preparing them according to training data) and apply gates only to hidden qubits.
Measure relevant observables: one can return probabilities (probs) or expectation values (expval) of Pauli operators, depending on the chosen loss function.
Define a cost function, such as the negative log-likelihood or a distance between output and target distribution.
Use an optimizer (e.g. gradient descent, Adam) with PennyLane’s gradient calculations to update parameters.




Below is a skeleton illustrating training a QCBM on a simple 2-point dataset. This is pseudocode for illustration only:

data = np.array([[0,1], [1,0]])  # two data points (|01> and |10>)

def cost(params):
    loss = 0
    for v in data:
        # clamp visible by preparing state |v>
        @qml.qnode(dev)
        def circuit_clamped():
            if v[0]==1: qml.PauliX(wires=0)
            if v[1]==1: qml.PauliX(wires=1)
            # apply layers on all qubits (visible hidden together in this toy example)
            qml.RY(params[0], wires=0)
            qml.RY(params[1], wires=1)
            return qml.probs(wires=[0,1])
        prob = circuit_clamped()
        # compute negative log prob of data state (this is a toy loss)
        idx = v[0]*2 + v[1]  # index of bitstring in probs
        loss += -np.log(prob[idx] + 1e-6)
    return loss / len(data)

opt = qml.GradientDescentOptimizer(stepsize=0.1)
for step in range(50):
    params, curr_cost = opt.step_and_cost(cost, params)
    if step%10==0:
        print(f"Step {step}: cost = {curr_cost:.3f}")
In this toy loop, we optimize the circuit so that the probability of the given data points is high. In a true QBM, the cost would be derived from the log-likelihood of all data under the quantum thermal model.



Exercises (Chapter 5): (1) PennyLane circuit: Modify the above QCBM code to use 3 qubits and train to model a three-bit data distribution. (2) Hamiltonian expectation: Using PennyLane, compute the expectation value of $H = \sigma^x\otimes\sigma^x + \sigma^z\otimes\sigma^z$ on the Bell state $(|00>+|11>)/\sqrt{2}$. (3) Parameter shift rule: Explain how PennyLane can compute gradients of quantum circuit parameters via the parameter-shift rule, and why this is useful for QBM training.





Applications and Examples





Although QBMs are still mainly a research topic, we briefly mention some toy applications and encourage experimentation with small datasets.  Classical RBMs have been used for simple generative tasks like the Bars-and-Stripes dataset (binary images) or the MNIST digit dataset.  QBMs could similarly be tested on these tasks to see if they offer any advantage. For example, one can train a small QBM to generate $4\times4$ binary patterns.  Tan et al. (2023) used QBMs on the Bars-and-Stripes task (4×4 images) as a benchmark .



In reinforcement learning, QBM-like models have been used to learn policies: one treats (state,action) pairs as inputs and uses the QBM to approximate a Q-function .  An RQBM with input neurons (for state), output neurons (for action), and hidden qubits can learn the value of each action by minimizing a “free energy” (analogous to Q-value).  This demonstrates how QBMs could serve as function approximators.



Small dataset example: Consider a 4-bit dataset consisting of the two strings 0000 and 1111. A classical RBM with sufficient hidden units can exactly model this. A QBM with 4 visible qubits and some hidden qubits could also learn it. One could set up the QBM Hamiltonian and run a variational algorithm so that $|0000>$ and $|1111>$ have highest probability.  Training could be done by maximizing $p(0000)+p(1111)$. This is an example of a bimodal distribution that a BM learns.



Another toy example is the XOR problem: classify inputs $[0,0],[0,1],[1,0],[1,1]$ with target labels given by XOR. A small QBM or RQBM could be used as a generative model over the four inputs and supervised labels, similar to a discriminative model. Although this is artificial, it illustrates how a hybrid quantum model might handle simple logical relationships.



Practical tips: When experimenting, it is often easiest to start with fully visible QBMs (no hidden units). In that case, the model distribution is $p(v) \propto \bra{v}e^{-\beta H}\ket{v}$ where $v$ runs over visible states. Gradient descent on this model can be done classically for very small sizes by explicitly computing $e^{-\beta H}$ and its gradients. For larger sizes, one must rely on quantum sampling or approximations.



Example Problem: Try training a 2-qubit QBM (4-dimensional state space) to match a given target distribution, e.g.\ $p(00)=0.5, p(11)=0.5$ and $p(01)=p(10)=0$. This requires choosing a Hamiltonian $H = a,\sigma^z_1\sigma^z_2 + b,(\sigma^z_1+\sigma^z_2)$ (and $\Gamma$ fields). Intuitively, $a$ should be negative (to favor $00,11$) and $b$ controls bias. Use gradient descent to find $a,b$ that achieve the distribution at some $\beta$.





References





Montúfar, G. (2018). Restricted Boltzmann Machines: Introduction and Review. Journal of Machine Learning Research 23(36), 1-108 .
Amin, M. H., Andriyash, E., Rolfe, J., Kulchytskyy, B., & Melko, R. (2018). Quantum Boltzmann Machine. Physical Review X, 8(2), 021050 .
Askarzadeh, A., Cong, I., & Sarma, S. (2024). On the sample complexity of quantum Boltzmann machine learning. Communications Physics 7, 24 .
Wiebe, N., & Wossnig, L. (2019). Generative training of quantum Boltzmann machines with hidden units. arXiv:1905.09902 .
Tan, A.T.V., Bahr, K., Plewnia, R., Grimsley, H.R., & Cerezo, M. (2023). Training quantum Boltzmann machines with coresets. arXiv:2307.14459 .
Hinton, G. E. (2002). Training products of experts by minimizing contrastive divergence. Neural Computation 14, 1771–1800.
Hinton, G. E., Osindero, S., & Teh, Y.-W. (2006). A Fast Learning Algorithm for Deep Belief Nets. Neural Computation 18(7), 1527–1554.






Index





Boltzmann machine (BM): Pairwise energy-based model; probabilistic graphical model (Ch. 2).
Partition function: Normalization constant $Z = \sum_x e^{-E(x)}$ in BM probability (Ch. 2).
Restricted Boltzmann machine (RBM): Bipartite BM with visible–hidden layers only (Ch. 2).
Quantum Boltzmann machine (QBM): Quantum neural network with qubit-based energy (Ch. 3, 4).
Restricted QBM (RQBM): Bipartite QBM with visible and hidden qubits (Ch. 4).
Gibbs state: Thermal state $\rho=e^{-\beta H}/\Tr(e^{-\beta H})$ in a quantum model (Ch. 3).
Relative entropy: Quantum KL divergence used as training objective (Ch. 4).
Contrastive Divergence (CD): Approximate BM training algorithm (Ch. 2, 6).
PennyLane: Quantum software library for circuit simulations (Ch. 5).
Hamiltonian: Quantum operator representing energy (Ch. 3, 5).
Visibility (clamping): Fixing visible qubits to data values during training (Ch. 3–5).
